% Results — tables in sections/generated/ sourced from results/tables/

Four side-effect dimensions meet pre-specified primary-claim thresholds on the Pass~1-flagged subset ($N{=}16{,}015$ parseable Pass~2 rows): steering asymmetry, inverse logit polarity, token inflation, and factual reversal.
Unless noted, Pass~2 statistics describe effects conditional on Pass~1 flagging (55.6\% of the 28{,}807-sample corpus), not corpus-wide prevalence.
Table~\ref{tab:primary_claims_summary} summarizes severity and prevalence for these four dimensions. Table~\ref{tab:dimension_claims} classifies all 19 dimensions.
Figure~\ref{fig:per_alpha_objective} summarizes multiplier dependence. Appendix~\ref{ap:iaa} reports dataset$\times$dimension severity heterogeneity (Figure~\ref{fig:heatmap}).

\begin{table}[htbp]
  \centering\small
  \caption{Primary-claim dimensions on the Pass~1-flagged subset ($N{=}16{,}015$). Prev.\ is pooled over Pass~2 rows unless noted as dataset-mean. Mean severities, Cohen's $d$, BH-adjusted tests, and dataset-level dispersion appear in Section~\ref{sec:results} and Appendix~\ref{ap:iaa}.}
  \label{tab:primary_claims_summary}
  \resizebox{\linewidth}{!}{%
  \begin{tabular}{@{}lrrrl@{}}
    \toprule
    Dimension & Mean sev. & Cohen's $d$ & Prev. & Primary cells \\
    \midrule
    Steering asymmetry & 2.29 & 1.61 & 75.9\% pooled & 36/36 \\
    Inverse logit polarity & 1.04 & 0.63 & 38.0\% dataset-mean & 29/36 \\
    Token inflation & 1.03 & 0.85 & 32.6\% dataset-mean & 32/36 \\
    Factual reversal & 1.02 & 0.53 & 20.6\% pooled / 23.6\% dataset-mean & 16/36 \\
    \bottomrule
  \end{tabular}%
  }
\end{table}


\subsection{Conditional Evaluation Sample}\label{Conditional Evaluation Sample}
Pass~1 screened every sample for steerability and broad anomaly categories, flagging 16{,}021 samples (55.6\%) for Pass~2; six flagged rows lacked parseable Pass~2 output (Table~\ref{tab:pass1_pass2_funnel}).
Pass~1 steerability was predominantly weak or absent: 51.9\% \textsc{Weak}, 21.9\% \textsc{None}, 12.5\% \textsc{Moderate}, 13.4\% \textsc{Strong}, and 0.2\% \textsc{Anti-Steered}.
Flag rates ranged from 23.9\% to 100.0\% across datasets (Appendix~\ref{ap:iaa}, Figure~\ref{fig:pass1_flags}).

\begin{table}[htbp]
    \centering
    \caption{Pass~1$\to$Pass~2 screening funnel (full-scale evaluation).}
    \label{tab:pass1_pass2_funnel}
    \begin{tabular}{lr}
        \toprule
        Stage & $N$ (or \%) \\
        \midrule
        Total corpus & 28{,}807 \\
        Pass-1 any category flagged & 16{,}021 (55.6\%) \\
        Pass-2 evaluated (Pass-1-flagged subset) & 16{,}021 \\
        Pass-2 any severity $> 0$ & 14{,}372 (89.7\% of Pass-2) \\
        Pass-2 all severities $= 0$ & 1{,}649 (10.3\% of Pass-2) \\
        \bottomrule
    \end{tabular}
\end{table}


\subsection{Primary Side Effects}
\textbf{Steering asymmetry} is the strongest and most universal primary claim: all 36 datasets are BH-significant and meet primary thresholds (Table~\ref{tab:dimension_claims}), with the largest mean severity and Cohen's $d$ (Table~\ref{tab:primary_claims_summary}).

\textbf{Token inflation}, \textbf{inverse logit polarity}, and \textbf{factual reversal} also carry primary claims but with greater dataset heterogeneity.
Inverse logit polarity has the highest inter-rater agreement among primary dimensions ($\alpha_{\mathrm{ord}}{=}0.87$) and non-zero judge ratings align with polarity inferred directly from target-token logits (precision 1.0, recall 0.99) (Appendix~\ref{ap:iaa}).
Human-LLM concordance on the four primary dimensions is reported in Appendix~\ref{ap:human_validation}.

Persona amplification and tonal rigidity are frequently BH-significant yet fail the IAA gate for primary claims.
Hedging escalation fails both the IAA and mean-severity thresholds.
\textbf{Logit-text decoupling} remains rare and exploratory.

Case studies can be found in Appendix~\ref{ap:case_studies}.

\begin{table}[htbp]
  \centering\small
  \caption{Primary vs.\ exploratory claims by side-effect dimension (36 datasets each). Primary: BH-significant with mean $\ge 0.5$, prevalence $\ge 10\%$, and IAA $\alpha_{\mathrm{ord}} \ge 0.65$.}
  \label{tab:dimension_claims}
  \begin{tabular}{@{}lrrrr@{}}
    \toprule
    Dimension & BH-sig. & Primary & Explor. & Mean sev. \\
    \midrule
    Steering asymmetry & 36 & 36 & 0 & 2.29 \\
    Token inflation & 36 & 32 & 4 & 1.03 \\
    Inverse logit polarity & 36 & 29 & 7 & 1.04 \\
    Factual reversal & 31 & 16 & 15 & 1.02 \\
    Option-label mismatch & 14 & 10 & 4 & 0.26 \\
    Factual drift & 11 & 4 & 7 & 0.19 \\
    Repetition looping & 25 & 4 & 21 & 0.36 \\
    Label-content contradiction & 13 & 4 & 9 & 0.18 \\
    Sensitivity refusal & 35 & 4 & 31 & 0.19 \\
    Hallucinated specificity & 25 & 2 & 23 & 0.16 \\
    Persona amplification & 35 & 0 & 35 & 0.69 \\
    Hedging escalation & 34 & 0 & 34 & 0.35 \\
    Token compression & 35 & 0 & 35 & 0.26 \\
    Neutralization rephrasing & 33 & 0 & 33 & 0.30 \\
    Implicit assumption injection & 29 & 0 & 29 & 0.12 \\
    Safety override resistance & 15 & 0 & 15 & 0.15 \\
    Monotonic factual scaling & 5 & 0 & 5 & 0.03 \\
    Logit-text decoupling & 19 & 0 & 19 & 0.12 \\
    Tonal rigidity & 36 & 0 & 36 & 0.64 \\
    \bottomrule
  \end{tabular}
\end{table}


\subsection{Statistical Power and Practical Significance}\label{Statistical Power and Practical Significance}
We test whether Pass~2 severity differs from zero using one-sided Wilcoxon signed-rank tests with Benjamini-Hochberg correction across 684 dataset-dimension pairs (Section~\ref{sec:Statistical Analysis}).
With $N \approx 500$-820 flagged samples per dataset, \textbf{503/684 tests (73.5\%) are BH-significant}, combining large prevalent effects with sparse ones driven by sample size rather than magnitude.
Pre-specified primary-claim criteria reduce this to \textbf{141 primary} and \textbf{362 exploratory} cells.
Table~\ref{tab:practical_sig} shows sensitivity to alternative cutoffs.
Pooled and dataset-mean prevalence can differ because the former sums non-zero rows over Pass~2 scored rows while the latter averages per-dataset rates.

\begin{table}[htbp]
    \centering
    \caption{Practical-significance sensitivity: BH-significant tests meeting mean-severity and prevalence thresholds.}
    \label{tab:practical_sig}
    \begin{tabular}{rrr}
        \toprule
        Mean $\ge$ & Prevalence $\ge$ (\%) & $N$ cells \\
        \midrule
        0.2 & 1 & 322 \\
        0.5 & 10 & 212 \\
        1.0 & 10 & 107 \\
        \bottomrule
    \end{tabular}
\end{table}


\subsection{Severity Across Steering Multipliers}\label{permultiplierresult}
Pass~2 overall severity scores assign one rating per dimension per sample after viewing all nine multipliers.
To test whether effects concentrate at specific strengths, we complement these overall scores with automated per-multiplier proxies: token-length ratio, text deviation from baseline, and mean $|\Delta\mathrm{logit}|$ relative to $\alpha{=}0$.

All three proxies grow monotonically with $|\alpha|$ on the flagged subset (Figure~\ref{fig:per_alpha_objective}), supporting the token-inflation primary claim independently of judge subjectivity.

Corpus-mean generation perplexity changes only modestly across the multiplier grid: $+5.6\%$ at $\alpha{=}{-}2.0$ and $+1.0\%$ at $\alpha{=}{+}2.0$ relative to $\alpha{=}0$ on the full corpus (Figure~\ref{fig:perplexity_strength}).
Perplexity can therefore remain stable while side effects intensify.
On the Pass~2-evaluated subset ($N{=}16{,}015$), per-sample max $|\Delta\mathrm{perplexity}|$ correlates weakly with overall severity on three primary dimensions (Pearson $\rho{=}0.36$ token inflation, $\rho{=}0.30$ factual reversal, $\rho{=}0.26$ steering asymmetry) and is essentially uncorrelated with inverse logit polarity ($\rho{\approx}{-}0.05$, Figure~\ref{fig:perplexity_side_effects}, Appendix~\ref{ap:iaa}).

A per-$\alpha$ judge validation subsample ($N{=}400$, $n{=}262$ comparable triples) corroborates overall Pass~2 scores on five dimensions (Appendix~\ref{ap:iaa}).

\begin{figure}[htbp]
    \centering
    \includegraphics[width=\textwidth]{per_alpha_objective_metrics.png}
    \caption{Automated per-multiplier proxies on the Pass-1-flagged subset ($N{=}16{,}015$). All three panels share steering multiplier $\alpha$ (x-axis). Left to right: mean token-length ratio vs.\ baseline, mean text deviation, and mean $|\Delta\mathrm{logit}|$ vs.\ $\alpha{=}0$; error bars show standard error of the mean.}
    \label{fig:per_alpha_objective}
\end{figure}

\begin{figure}[htbp]
    \centering
    \includegraphics[width=0.75\textwidth]{perplexity_vs_steering_strength.png}
    \caption{Mean generation perplexity across $\alpha$ on the full corpus ($N{=}28{,}807$, blue) and Pass~2-evaluated subset ($N{=}16{,}015$, orange). Error bars show standard error of the mean.}
    \label{fig:perplexity_strength}
\end{figure}
